\exercisegroup

\begin{enumerate}
\def\labelenumi{\arabic{enumi})}
\item
  设 A 是一个凸集。那么，一点 \(x \in A\) 是 A 的极点，当且仅当对 A 的任何子集 B，由 x 属于 B 的凸包这一论断可推出 \(x \in B\) 。
\item
  沿用 II, 第 74 页, 习题 3 的记号，设 G 是 E 中由 K ∪ \{λ\} 生成的向量子空间。证明：在 G 中，点 λ 是 K 的闭凸包的一个极点，但 λ 不属于 K（参见 II, 第 25 页, 推论）。
\end{enumerate}

¶ 3) 设 A 是向量空间 E 中的一个凸集，x 是 A 的一点。我们把由 x 以及 A 中这样的 \(y \neq x\) （过 x 与 y 的直线包含一个开线段，它含于 A 且含 x）所构成的集合称为 x 在 A 中的刻面。A 的（关于由 A 生成的线性流形的）代数内点（相应地，极点）（II, 第 26 页）是指在 A 中的刻面等于 A（相应地，缩为一点）的那些点。

\begin{enumerate}
\def\labelenumi{\alph{enumi})}
\item
  证明：刻面 \(\mathrm{F}_x\) （一点 \(x \in \mathbf{A}\) 的刻面）是最大的凸集 \(\mathbf{B} \subset \mathbf{A}\) ，使得 \(x\) 是 \(\mathbf{B}\) （关于由 \(\mathbf{B}\) 生成的线性流形）的代数内点。
\item
  对每一点 \(y \in \mathbf{F}_x\) ，刻面 \(\mathbf{F}_y\) （\(y\) 在 \(\mathbf{A}\) 中的刻面）与 \(y\) 在 \(\mathbf{F}_x\) 中的刻面恒同。要使 \(\mathbf{F}_y = \mathbf{F}_x\) ，必须且只需 \(y\) 是 \(\mathbf{F}_x\) （关于由 \(\mathbf{F}_x\) 生成的线性流形）的代数内点。由此推出：如果 \(\mathbf{F}_x\) 是有限维的，且 \(y\) 是 \(\mathbf{F}_x\) （关于由 \(\mathbf{F}_x\) 生成的线性流形）的非代数内点，那么 \(\mathbf{F}_y\) 的维数严格小于 \(\mathbf{F}_x\) 的维数。
\item
  E 中与 A 相交且满足下述条件的线性流形 V 称为 A 的支撑流形：对每个 \(x \in A \cap V\) ，每个含于 A 且含 x 的开线段必然含于 V。证明：对所有 \(x \in A\) ，由 x 在 A 中的刻面 \(F_{x}\) 生成的线性流形 M 是包含 x 的最小支撑流形，且 \(M \cap A = F_{x}\) 。对 A 的每个支撑流形 V，交 \(V \cap A\) 是它的每个代数内点（关于由 \(V \cap A\) 生成的线性流形）在 A 中的刻面。
\end{enumerate}

\(d\) ) 设 \(\mathbf{A}\) 与 \(\mathbf{B}\) 是 \(\mathbf{E}\) 中两个凸集。对每一点 \(x\in \mathbf{A}\cap \mathbf{B}\) ，\(x\) 在 \(\mathbf{A}\cap \mathbf{B}\) 中的刻面是 \(x\) 在 \(\mathbf{A}\) 中与在 \(\mathbf{B}\) 中诸刻面之交。

\begin{enumerate}
\def\labelenumi{\alph{enumi})}
\setcounter{enumi}{4}
\item
  设 B 是 Hausdorff 拓扑向量空间 E 中的一个闭凸集，且 B 含有一个具有有限余维数 n 的闭线性流形 M；那么 B 的每一点在 B 中的刻面都含有一个余维数为 n 的闭线性流形（II, 第 67 页, 习题 14, d)）。如果 A 是一个凸集，那么 A ∩ B 的一点 x 在 A ∩ B 中的刻面是有限维的，当且仅当 x 在 A 中的刻面是有限维的：进而，若 p 与 q 分别是 x 在 A 中的刻面与 x 在 A ∩ B 中的刻面的维数，则 \(p \leqslant q + n\) 。特别地，如果 \(x \in A \cap B\) 是 A ∩ B 的一个极点，那么它在 A 中的刻面的维数 \(\leqslant n\) 。
\item
  由 e) 推出：如果 A 是紧的，且 V 是 E 中一个具有有限余维数 n 的闭线性流形，那么 \(V \cap A\) 的每个极点至多是 A 的 \(n + 1\) 个极点的线性组合。
\end{enumerate}

\begin{enumerate}
\def\labelenumi{\arabic{enumi})}
\setcounter{enumi}{3}
\item
  在平面 \(\mathbf{R}^2\) 中，考虑凸集 \(\mathbf{A}\) ，它由这样的点 \((\xi, \eta)\) 构成：\(-1 \leqslant \xi \leqslant 1\) ，\(-1 - \sqrt{1 - \xi^2} \leqslant \eta \leqslant 1 + \sqrt{1 - \xi^2}\) 。证明：存在 \(\mathbf{A}\) 的一些边界点，它们的刻面异于 \(\mathbf{A}\) 与过该点的 \(\mathbf{A}\) 的支撑直线之交。
\item
  在实数的有界序列所成的 Banach 空间 \(l^{\infty}(\mathbf{N})\) （其中的序列写作 \(x = (\xi_n)\) ）中，设 \(\mathbf{A}\) 是由不等式 \(-1 / n \leqslant \xi_n \leqslant 1\) （对 \(n \geqslant 1\) ）与 \(-1 \leqslant \xi_0 \leqslant 1\) 所定义的闭凸集。证明：\(\mathbf{A}\) 有非空内部，原点是 \(\mathbf{A}\) 的一个边界点，且 0 在 \(\mathbf{A}\) 中的刻面不是闭的。如果我们给 \(\mathbf{A}\) 装备由乘积空间 \(\mathbf{R}^{\mathbf{N}}\) 的拓扑诱导的拓扑，证明 \(\mathbf{A}\) 是紧的，但 0 在 \(\mathbf{A}\) 中的刻面在 \(\mathbf{A}\) 中不闭。
\item
  设 E 与 \(E'\) 是处于分离对偶中的两个向量空间，A 是 E 中包含 0 且对 \(\sigma(\mathrm{E}, \mathrm{E}')\) 闭的一个凸集。对所有 \(a \in A\) ，记 \(F_{a}'\) 为点 \(x' \in A^{\circ}\) （满足 \(\langle a, x' \rangle = -1\) ）所成之集；它是一个闭的（对 \(\sigma(\mathrm{E}', \mathrm{E})\) ）凸集，含于 \(A^{\circ}\) 。证明：\(F_{a}'\) 是 \(A^{\circ}\) 中的刻面，即 \(F_{a}'\) （关于由 \(F_{a}'\) 生成的线性流形）的每个代数内点的刻面。我们称 \(F_{a}'\) 为 a 在 \(A^{\circ}\) 中的对偶刻面。如果 \(F_{a}\) 是 a 在 A 中的刻面，证明：\(F_{a}'\) 也是在 \(A^{\circ}\) 中的对偶刻面——对 \(F_{a}\) （关于由 \(F_{a}\) 生成的线性流形）的每个代数内点而言；进而，如果把 A 等同于 \(A^{\circ\circ}\) ，那么 \(F_{a}^{\prime}\) （关于由 \(F_{a}^{\prime}\) 生成的线性流形）的每个代数内点在 A 中的对偶刻面包含 \(F_{a}\) 。当 \(F_{a}^{\prime}\) 非空时（当 E 有限维且 \(a \neq 0\) 时这总是成立的，参见 II, 第 78 页, 习题 13），我们称 \(F_{a}\) 与 \(F_{a}^{\prime}\) 中的每一个是另一个的对偶刻面。
\end{enumerate}

我们称点 \(a \in \mathbf{A}\) 是 \(\mathbf{A}\) 的光滑点，如果 \(\mathbf{F}_a'\) 缩为一点（换句话说，如果存在 \(\mathbf{A}\) 的一个过 \(a\) 的闭支撑超平面）；我们称 \(a\) 是严格凸点（或暴露点），如果存在一个闭支撑超平面 \(\mathbf{H}\) （它支撑 \(\mathbf{A}\) ），使得 \(\mathbf{H} \cap \mathbf{A} = \{a\}\) ；这等价于说：\(\mathbf{F}_a'\) （关于由 \(\mathbf{F}_a'\) 生成的某个线性流形）的一个代数内点是 \(\mathbf{A}^\circ\) 的光滑点。

\begin{enumerate}
\def\labelenumi{\arabic{enumi})}
\setcounter{enumi}{6}
\tightlist
\item
  设 \(\mathbf{E}\) 是有限维（维数 \(n\) ）的向量空间，\(\mathbf{A}\) 是 \(\mathbf{E}\) 中的一个闭凸集，0 是它的一个内点。
\end{enumerate}

\begin{enumerate}
\def\labelenumi{\alph{enumi})}
\item
  设 F 与 \(F'\) 是 A 与 \(A^{\circ}\) 的两个对偶刻面（习题 6）；如果 F 是 p 维的而 \(F'\) 是 q 维的，试证明 \(p + q \leqslant n - 1\) 。对 A 的每个边界点 x，x 在 A 中的刻面的维数称为 x 的阶，其对偶刻面在 \(A^{\circ}\) 中的维数称为 x 的类。A 的刻面 F 的阶（相应地，类）按定义是 F 的某个代数内点（关于由 F 生成的线性流形）的阶（相应地，类）。A 的极点是 0 阶点；A 的光滑点（习题 6）是 0 类点。
\item
  A 的一个 n - 1 类（从而 0 阶）的边界点称为 A 的顶点。证明：A 的顶点所成之集是可数的（考察 A 的诸顶点的对偶刻面所成之集，GT, VI, § 2, 习题 12）。
\item
  设 F 是 A 的一个 p 维刻面，M 是 n - p 维的一个线性流形，它与 F 仅交于一点 a，a 是 F 的一个代数内点，且 M 含有 A 的一个内点。证明：如果 V 是 \(M \cap A\) 在 M 中的一个过 a 的支撑超平面，那么由 \(F \cup V\) （在 E 中）生成的超平面 H 是 A 的一个支撑超平面。
\item
  我们称 A 的一个 p 阶 q 类的刻面 F 为超刻面，如果 \(p + q = n - 1\) ；此时其对偶刻面也是 \(A^{\circ}\) 的一个超刻面。如果一个 n - p 维的线性流形 M 与一个超刻面 F 仅交于一点，且该点是 F 的一个代数内点（关于由 F 生成的线性流形），证明：该点是凸集 \(M \cap A\) 的一个顶点，反之亦然（利用 c)）。由此推出：A 的 p 阶超刻面所成之集是可数的。（把 E 等同于 \(R^{n}\) ，考察 A 在 \(R^{n}\) 的每个 p 维坐标流形上的投影；如果其在一个坐标流形 V 上的投影为 p 维的那种 p 阶超刻面所成之集不可数，则借助 V 中具有有理坐标的点，证明存在 V 的一点，它是不可数无穷多个这种投影的内点；然后利用 b)。）试给出一个凸集的例子，它有不可数无穷多个刻面，其中每一个既不是单独一点也不是超刻面。
\item
  如果 A 的所有边界点都是光滑的，证明：把 A 的边界 G 的每一点 x 对应到 x 的对偶刻面的唯一点的映射，是 G 到 A° 的边界上的一个连续映射（参见 TG, I, § 9.1, 推论）。在什么情形下这个映射是双射？
\end{enumerate}

\begin{enumerate}
\def\labelenumi{\arabic{enumi})}
\setcounter{enumi}{7}
\tightlist
\item
  设 \(\mathbf{E}\) 是有限维（维数 \(n\) ）的向量空间，\(\mathbf{A}\) 是 \(\mathbf{E}\) 中的一个紧凸集。
\end{enumerate}

\begin{enumerate}
\def\labelenumi{\alph{enumi})}
\item
  设 H 是 E 中一个超平面。证明：在由 H 确定的、且至少含 A 的一点的开半空间中，存在 A 的一个严格凸点（II, 第 87 页, 习题 6）。（在 H 中考虑一个 n - 1 维的、半径充分大且包含 H ∩ A 的闭欧氏球 C，然后考虑包含 A 且满足 B ∩ H = C 的、半径更大的 n 维欧氏球 B。）
\item
  证明：A 是严格凸点所成之集的闭凸包（利用 a)）。c) 证明：A 的每个极点都是 A 的严格凸点所成之集的一个聚点。（利用 b) 与 II, 第 66 页, 习题 9, a)，注意一个极点是形如 \(\sum_{i=0}^{n} \lambda_{im} x_{im}\) 的点组成的序列的极限，其中 \(\lambda_{im} \geqslant 0\) ，\(\sum_{i=0}^{n} \lambda_{im} = 1\) ，且诸 \(x_{im}\) 是 A 的严格凸点；然后利用 A 的紧性。）
\end{enumerate}

\begin{enumerate}
\def\labelenumi{\arabic{enumi})}
\setcounter{enumi}{8}
\item
  证明：在乘积空间 \(\mathbf{E} = \mathbf{R}^{\mathbf{N}}\) 中，方体 \(\mathbf{I}^{\mathbf{N}}\) （其中 \(\mathbf{I} = [0,1]\) ）是没有严格凸点的紧凸集。
\item
  在空间 \(\mathbf{R}^2\) 中，证明：闭凸集 \(\mathbf{A}\) 的极点所成之集是闭的（证明 A 中其在 A 中的刻面为 1 维的那种点所成之集，相对于 A 的边界是开集）。
\item
  \begin{enumerate}
  \def\labelenumii{\alph{enumii})}
  \tightlist
  \item
    在空间 \(\mathbf{R}^3\) 中，考虑紧凸集 A，它是圆 \(\zeta = 0\) ，\(\xi^2 + \eta^2 - 2\xi = 0\) 与两点 \((0, 0, 1)\) 、\((0, 0, -1)\) 之并的凸包。证明：A 的极点所成之集在 A 中不闭。
  \end{enumerate}
\end{enumerate}

\begin{enumerate}
\def\labelenumi{\alph{enumi})}
\setcounter{enumi}{1}
\tightlist
\item
  设 A 是 Hausdorff 拓扑向量空间 E 中一个可度量化的紧凸集。证明：A 的极点所成之集是 A 中一列开集之交。（如果 d 是定义 A 的拓扑的一个距离，那么对每个整数 n，考虑点 \(x = \frac{1}{2}(y + z)\) 所成之集，其中 y, z 属于 A 且 \(d(y, z) \geqslant 1/n\) 。）
\end{enumerate}

\begin{enumerate}
\def\labelenumi{\arabic{enumi})}
\setcounter{enumi}{11}
\tightlist
\item
  在 Banach 空间 \(l^{\infty}(\mathbf{N})\) 中，设 \(e_n\) 是除第 \(n\) 项为 1 外其余各项均为零的序列。设 \(\mathbf{A}\) 是由 0 与诸点 \(\mathbf{e}_n / (n + 1)\) ( \(n \geqslant 0\) )所成之集的闭凸包。证明：\(\mathbf{A}\) 是紧的，但它不等同于其极点所成之集的凸包。
\end{enumerate}

* 13) 在 Hilbert 空间 \(l^2(\mathbf{N})\) 中，设 \(\mathbf{A}\) 是这样的点 \(x = (\xi_n)\) 所成之集：\(\sum_{n} 2^{2n}\xi_n^2 \leqslant 1\) 。证明：\(\mathbf{A}\) 是凸的、紧的，且它是其极点所成之集的闭包。

\begin{enumerate}
\def\labelenumi{\arabic{enumi})}
\setcounter{enumi}{13}
\tightlist
\item
  设 \(\mathbf{E}\) 是 Banach 空间 \(l^{\infty}(\mathbf{N})\) 的一个闭向量子空间，它由序列 \(x = (\xi_n)\) （满足 \(\lim_{n\to \infty}\xi_n = 0\) ）组成。
\end{enumerate}

\begin{enumerate}
\def\labelenumi{\alph{enumi})}
\tightlist
\item
  证明：在 Banach 空间 E 中，闭单位球 B 没有任何极点。b) 设 \(u\) 是 E 上的连续线性型 \((\xi_n) \mapsto \sum_n 2^{-n}\xi_n\) 。证明：不存在 B 的支撑超平面平行于方程为 \(u(x)=0\) 的闭超平面。
\end{enumerate}

\begin{enumerate}
\def\labelenumi{\arabic{enumi})}
\setcounter{enumi}{14}
\tightlist
\item
  设 \(\mathbf{A}\) 是赋范空间 \(\mathbf{E}\) 中的一个紧集。
\end{enumerate}

\begin{enumerate}
\def\labelenumi{\alph{enumi})}
\item
  证明：A 的两个平行支撑超平面之间的距离至多等于 A 的直径 \(\delta\) 。
\item
  证明：存在 A 的点偶 \((a, b)\) 使 \(\|a - b\| = \delta\) ；对这样的点偶，存在 A 的两个平行支撑超平面，分别过 a 与 b，且它们之间的距离为 \(\delta\) （考察中心在 a、半径为 \(\delta\) 的闭球）。
\end{enumerate}

\begin{enumerate}
\def\labelenumi{\arabic{enumi})}
\setcounter{enumi}{15}
\tightlist
\item
  \begin{enumerate}
  \def\labelenumii{\alph{enumii})}
  \tightlist
  \item
    设 A 是空间 \(R^{n}\) （带欧氏范数）中一个 n 维紧凸集；对每个 \(z \in S_{n-1}\) ，用 \(\rho(z)\) 表示平行于向量 z 且含于 A 的线段长度的上确界。证明：存在 A 的两点 u, v，使得端点为 u, v 的线段平行于 z 且长度为 \(\rho(z)\) ；由此推出：存在 A 的两个平行的支撑超平面，分别过 u 与 v（考察集合 \(A' = A + \rho(z)z\) ，并用一个超平面分离集合 A 与 \(A'\) ）。
  \end{enumerate}
\end{enumerate}

\begin{enumerate}
\def\labelenumi{\alph{enumi})}
\setcounter{enumi}{1}
\tightlist
\item
  设 d 是 A 的两个平行支撑超平面之间距离的下确界；证明：存在 A 的两点 a, b 使 \(\|a - b\| = d\) ，且分别过 a 与 b 并垂直于 a - b 的超平面都是 A 的支撑超平面（利用 a)）。
\end{enumerate}

\begin{enumerate}
\def\labelenumi{\arabic{enumi})}
\setcounter{enumi}{16}
\item
  在空间 \(l^{\infty}(\mathbf{N})\) 中，设 \(\mathbf{A}\) 是 II, 第 89 页, 习题 12 中定义的紧凸集，并设 \(\mathbf{E}\) 是 \(l^{\infty}(\mathbf{N})\) 中由 \(\mathbf{A}\) 生成的闭向量子空间。证明：\(\mathbf{A}\) 在空间 \(\mathbf{E}\) 中的两个平行闭支撑超平面之间距离的下确界等于 0，尽管 \(\mathbf{A}\) 并不含于 \(\mathbf{E}\) 的任何闭超平面中。
\item
  在 Hausdorff 局部凸空间 E 中，设 \((\mathbf{K}_{\alpha})_{\alpha\in\mathbf{I}}\) 是由非空紧凸集组成的一个递减定向族。对所有 \(\alpha\in I\) ，把 \(K_{\alpha}\) 的极点所成之集记作 \(A_{\alpha}\) ，并把 \(F_{\alpha}\) 记作诸 \(A_{\beta}\) （\(\beta\geqslant\alpha\) ）之并的闭包，于是 \((\mathrm{F}_{\alpha})\) 是紧集的一个递减定向族。设 A 是诸 \(F_{\alpha}\) 的（非空）交，K 是诸 \(K_{\alpha}\) 的（非空）交。证明：K 是 A 的闭凸包。（如果 f 是 E 上的一个连续线性型，而 \(x_{\alpha}\) 是 \(F_{\alpha}\) 中使 f 在 \(F_{\alpha}\) 上达到其最大值的一点，证明：\(f(y)\leqslant f(x_{\alpha})\) 对所有 \(y\in K\) 成立；然后沿 I 的截线滤子取族 \((x_{\alpha})\) 的一个聚点。）
\end{enumerate}

¶ 19) 在空间 \(\mathbf{R}^n\) 中，设 \((\mathrm{K}_{\alpha})\) 是一族紧集，其个数 \(\geqslant n + 1\) ，且其中没有一个含于一个仿射超平面。假设对仿射自同构的每一族 \((u_{\alpha})\) （\(\mathbf{R}^n\) 的仿射自同构），如果任何 \(n + 1\) 个集合 \(u_{\alpha}(\mathrm{K}_{\alpha})\) 有公共点，那么所有 \(u_{\alpha}(\mathrm{K}_{\alpha})\) 有公共点。证明：在这些条件下诸 \(\mathbf{K}_{\alpha}\) 是凸的。（反之，设存在 \(n + 1\) 个点 \(x_1, \ldots, x_{n+1}\) 属于同一集合 \(\mathbf{K}_{\alpha}\) ，且存在一点 \(x_0\) 属于诸 \(x_i\) ( \(i \geqslant 1\) )所成之集的凸包，但不属于 \(\mathbf{K}_{\alpha}\) 。注意对每个指标 \(i \geqslant 1\) ，存在一个仿射自同构 \(u_i\) （\(\mathbf{R}^n\) 的）与一个指标 \(\alpha_i\) ，使得 \(x_0\) 与诸 \(x_j\) （指标 \(j \neq i\) 的那些）都是 \(u_i(\mathrm{K}_{\alpha_i})\) 的极点，并证明这 \(n + 2\) 个集合 \(\mathbf{K}_{\alpha}\) 与 \(u_i(\mathrm{K}_{\alpha_i})\) 没有公共点。）

\begin{enumerate}
\def\labelenumi{\arabic{enumi})}
\setcounter{enumi}{19}
\tightlist
\item
  试给出一个紧凸集 \(\mathbf{K}\) （在 \(\mathbf{R}^2\) 中）的例子，它包含 0，且以 0 为顶点、由 \(\mathbf{K}\) 生成的锥在 \(\mathbf{R}^2\) 中不闭。
\end{enumerate}

¶ 21) a) 在 Hausdorff 局部凸空间 E 中，设 A 是一个局部紧的闭凸锥，它不含任何直线。证明：A 是一个具有紧底的锥（对 A 的顶点——它是 A 的一个极点——应用 II, 第 55 页, 命题 2）。由此推出：存在 A 的一个闭支撑超平面 H，它含 A 的顶点 s 且 \(H \cap A = \{s\}\) 。b) 设 A, B 是 E 中两个以 0 为顶点的闭凸锥，它们都局部紧且不含直线。证明：如果 \(A \cap B = \{0\}\) ，那么 A - B 是一个闭的、局部紧的、不含任何直线的锥（方法同 II, 第 67 页, 习题 16）。由此推出：存在一个同时支撑 A 与 B 的闭超平面，它把 A 与 B 分离，且 \(H \cap A = H \cap B = \{0\}\) 。

\begin{enumerate}
\def\labelenumi{\alph{enumi})}
\setcounter{enumi}{2}
\tightlist
\item
  试给出一个局部紧闭凸锥 A 的例子，使 \(\mathbf{A} - \mathbf{A}\) 不是局部紧的（参见 II, 第 78 页, 习题 11）。
\end{enumerate}

\(d)\) 设 \(\mathbf{A}\) 是 \(\mathbf{E}\) 中一个局部紧的闭凸集，它不含直线。设 \(x_0\) 是 \(\mathbf{A}\) 的一点，\(\mathbf{C}_{\mathbf{A}}\) 是 \(\mathbf{A}\) 的渐近锥（参见 II, 第 67 页, 习题 14），\(\mathbf{H}\) 是 \(x_0 + \mathbf{C}_{\mathbf{A}}\) 的一个过 \(x_0\) 且满足 \((x_0 + \mathbf{C}_{\mathbf{A}}) \cap \mathbf{H} = \{x_0\}\) 的支撑超平面。如果 \(f(x) = a\) 是 \(\mathbf{H}\) 的方程，且 \(f(x) \geqslant a\) 在 \(x_0 + \mathbf{C}_{\mathbf{A}}\) 中成立，那么证明：对每个实数 \(b\) ，由 \(y \in \mathbf{A}\) （满足 \(f(y) \leqslant b\) ）所成之集是紧的。

¶ 22) 向量空间 E 的凸子集 A 的极射线是指一条含于 A 的闭半直线 D，使得：对所有 \(x \in \mathbf{D}\) 以及每个端点 \(a, b\) 在 A 中且包含 \(x\) 的开线段，必然有 \(a \in \mathbf{D}\) 且 \(b \in \mathbf{D}\) ；D 的端点是 A 的一个极点。

\begin{enumerate}
\def\labelenumi{\alph{enumi})}
\item
  在 Hausdorff 局部凸空间 E 中，证明：每个局部紧的、不含直线的闭凸集是其极射线与其极点之并的闭凸包。（设若不然，用 B 表示这个闭凸包，首先注意，利用习题 21, d)，存在一个闭超平面 H，使得 H ∩ A 紧且非空，而 H ∩ B = ∅。然后证明：如果 \(a \in H \cap A\) 是 H ∩ A 的一个极点（由假设它因而不是 A 的极点），且如果端点为 b, c、含于 A 而不含于 H 的开线段 S 包含 a，那么包含 S 的直线 D 必然包含一个含 a 且端点为 A 的极点的线段，或者包含 A 的一条含 a 的极射线。）
\item
  证明：如果 E 是有限维的，那么 E 中每个不含任何直线的闭凸集是其极点与极射线之并的凸包（对 E 的维数用归纳法论证）。
\end{enumerate}

\begin{enumerate}
\def\labelenumi{\arabic{enumi})}
\setcounter{enumi}{22}
\tightlist
\item
  在 \(R^{3}\) 中，考虑一个有内点的闭凸集 A，其边界 F 包含两个开线段 S, T，它们分别位于两条不平行的直线 D, \(D'\) 上（于是 S, T 的点在 A 中不是极点），而其所有其他边界点都是极点（可以说明如何定义这样的凸集）。对每个 \(x \in R^{3}\) ，令 \(f(x) = (d(x, \mathbf{D}))^{2}\) 。设 B 是 \(R^{4} = R^{3} \times R\) 中由偶 \((x, \zeta)\) （满足 \(x \in A\) 且 \(\zeta \geqslant f(x)\) ）所构成的闭凸集。证明：在 B 中，极点所成之集、极射线之并、极射线端点所成之集，以及这些集合中两个或三个的所有并，都是不闭且非空的。
\end{enumerate}

¶ 24) 在 \(\mathbf{R}^n\) 中，有限多个闭半空间（相应地，由过同一点的诸超平面所确定的闭半空间）的每个交称为多面体（相应地，多面锥）。凸集 \(\mathbf{C} \subset \mathbf{R}^n\) 在一点 \(x \in \mathbf{C}\) 处是局部多面体的，如果存在一个邻域 \(\mathbf{V}\) （\(x\) 在 \(\mathbf{C}\) 中的邻域），它是多面体。

\begin{enumerate}
\def\labelenumi{\alph{enumi})}
\item
  证明：如果闭凸集 \(\mathbf{C} \subset \mathbf{R}^n\) 在一点 \(x \in \mathbf{C}\) 处是局部多面体的，那么以 \(x\) 为顶点、由 \(\mathbf{C}\) 生成的锥是多面锥。
\item
  证明：\(\mathbf{R}^n\) 中在每个点处都局部多面体的紧凸集是多面体（利用 \(a\) ）。
\item
  设 \(\mathbf{P} \subset \mathbf{R}^n\) 是一个有内点的闭凸集。证明：下述条件等价：
\end{enumerate}

α) P 是多面体。

β) P 只有有限多个刻面（II, 第 87 页, 习题 3）。

\(\gamma)\) P 是一个集合的凸包，该集合是有限多个点与有限多条闭半直线之并。

（为证明 \(\alpha\) ) 蕴含 \(\beta\) ) ，把 P 取为数目尽可能少的闭半空间之交，并证明定义这些半空间的诸超平面由 P 的 \(n - 1\) 维刻面生成。为证明 \(\beta\) ) 蕴含 \(\gamma\) ) ，对 \(n\) 用归纳法论证。最后，为看出 \(\gamma\) ) 蕴含 \(\alpha\) ) ，考察 P 的极集 \(\mathrm{P}^{\circ}\) 。）

\begin{enumerate}
\def\labelenumi{\alph{enumi})}
\setcounter{enumi}{3}
\item
  证明：每个凸多面体 \(\mathbf{P}\) 都可以写成形式 \(\mathrm{Q} + \mathrm{C}_{\mathrm{P}}\) ，其中 \(\mathbf{Q}\) 是一个紧多面体，而 \(\mathrm{C}_{\mathrm{P}}\) 是 \(\mathbf{P}\) 的渐近锥。非紧的多面体不可能是抛物的。
\item
  证明：凸多面体的每个刻面都是超刻面（II, 第 88 页, 习题 7, d)）（对 n 用归纳法论证）。
\end{enumerate}

¶ 25) a) 设 \(\mathbf{C} \subset \mathbf{R}^n\) 是一个以 0 为顶点的闭凸锥。证明：C 在 \(\mathbf{R}^n\) 的每个 2 维子空间上的投影都是闭的，当且仅当 C 是多面锥（习题 24）。（化归为 C 不含直线的情形：利用 C 的紧底 S 的存在性（II, 第 90 页, 习题 21, a)），对 \(n\) 用归纳法论证，投影到与连接 0 和 S 的一个极点的直线平行的超平面上，并由此推出 S 是局部多面体的（习题 24）。b) 由 a) 推出：如果给 \(\mathbf{R}^n\) 装备使 C 成为 \(\geqslant 0\) 元素之集的序，那么 \(\mathbf{R}^n\) 的任何向量子空间 F 上的每个正线性型都可以延拓为 \(\mathbf{R}^n\) 上的正线性型，当且仅当 C 是多面锥（对极锥 \(\mathbf{C}^\circ\) 应用 a)）。如果 C 是 \(\mathbf{R}^3\) 中由 \((\xi_1, \xi_2, \xi_3)\) （满足 \(\xi_1 = 1\) ，\(\xi_3 \geqslant (\xi_2^3)^-\) ）生成的锥，而 F 是子空间 \(\xi_3 = 0\) ，试给出 F 上的一个正线性型的例子，它不能延拓为 \(\mathbf{R}^3\) 上的正线性型。

\begin{enumerate}
\def\labelenumi{\alph{enumi})}
\setcounter{enumi}{2}
\tightlist
\item
  设 A 是 \(\mathbf{R}^n\) 中一个多面体。那么，\(\mathbf{A} \cup \mathbf{B}\) 的凸包对每个多面体 B 都是闭的，当且仅当 A 是紧的（利用习题 24, d)）。
\end{enumerate}

\begin{enumerate}
\def\labelenumi{\arabic{enumi})}
\setcounter{enumi}{25}
\tightlist
\item
  \begin{enumerate}
  \def\labelenumii{\alph{enumii})}
  \tightlist
  \item
    设 E 是 Hausdorff 局部凸空间，A 是 E 中凸集 C 的一个帽。如果 \(s \in C\) 是不属于 A 的一点，且 B 是以 s 为顶点、由 A 生成的锥，证明 \(\mathbf{B} \cap (\mathbf{C} \cap \complement\mathbf{A})\) 的闭包是 B 中的一个帽。
  \end{enumerate}
\end{enumerate}

\begin{enumerate}
\def\labelenumi{\alph{enumi})}
\setcounter{enumi}{1}
\item
  假设 E 是有限维的。证明：E 中闭凸集 C 的每个帽 A 都可以用下述方式得到：考虑 C 的一个刻面 F（II, 第 87 页, 习题 3）以及由 F 生成的仿射线性流形 V 中的一个超平面 H，使 F 完全位于 H 的一侧，把 F 中位于 H 与 F 的一个平行于 H 的超平面 H' 之间的点所成之集取作 A（利用 a) 与 II, 第 38 页, 命题 4）。C 的刻面的每个极点是 C 的极点。
\item
  试在 Hausdorff 局部凸空间 E 中给出一个紧凸集 C 及其一个帽 A 的例子，使得 A 与 \(C \cap \complement A\) 各自都生成 E，且 A 与 \(C \cap \complement A\) 不能被 E 的一个闭超平面分离（参见 II, 第 78 页, 习题 11）。
\end{enumerate}

\begin{enumerate}
\def\labelenumi{\arabic{enumi})}
\setcounter{enumi}{26}
\item
  设 C 是直线的乘积 E 中一个闭凸集，\(a\) 是 C 的一个极点。证明：对 \(a\) 在 C 中的每个邻域 V，存在 E 中一个开半空间 F，使得 \(a \in \mathrm{F} \cap \mathrm{C} \subset \mathrm{V}\) 。（化归为 C 紧的情形。）
\item
  设 I 是一个不可数的无限集。证明：锥 \(R_{+}^{I}\) 在 \(R^{I}\) 中的每个帽都含于形如 \(R^{J}\) 的子空间的一个和式中，其中 J 是 I 的可数子集（利用 II, 第 38 页, 命题 4）。由此推出：存在 \(R_{+}^{I}\) 的点不属于 \(R_{+}^{I}\) 的任何帽，尽管 \(R_{+}^{I}\) 是其极生成元之并的闭凸包。
\item
  \begin{enumerate}
  \def\labelenumii{\alph{enumii})}
  \tightlist
  \item
    设 \((\mathrm{E}_n)\) 是 Hausdorff 局部凸空间的一个序列，\(\mathrm{E} = \prod_{n}\mathrm{E}_{n}\) 是它们的乘积。
  \end{enumerate}
\end{enumerate}

在每个 \(E_{n}\) 中，设 \(C_{n}\) 是以 0 为顶点的凸锥，\(A_{n}\) 是 \(C_{n}\) 的一个帽。证明：存在 \(C = \prod C_{n}\) 的一个帽，它包含 \(\prod A_{n}\) （按 II, 第 59 页, 命题 5 的方式论证）。

\begin{enumerate}
\def\labelenumi{\alph{enumi})}
\setcounter{enumi}{1}
\tightlist
\item
  设 \((\overline{\mathbf{E}}_n, \phi_{nm})\) 是 Hausdorff 局部凸空间的一个可数定向投影体系，\(\mathrm{E} = \varprojlim \mathrm{E}_n\) 是它的投影极限。对所有 \(n\) ，设 \(\mathbf{C}_n\) 是以 0 为顶点的凸锥，使得 \((\mathbf{C}_n)\) 是集合的一个投影体系。证明：如果对每个 \(n\) ，集合 \(C_n\) 都是它的诸帽之并，那么这对 \(\mathrm{C} = \varprojlim \mathrm{C}_n\) 也成立（利用 \(a\) ）。特别地，如果诸 \(\mathbf{C}_n\) 是具有紧底的锥，那么 \(\mathbf{C}\) 是其极生成元之并的闭凸包。
\end{enumerate}

\begin{enumerate}
\def\labelenumi{\arabic{enumi})}
\setcounter{enumi}{29}
\tightlist
\item
  设 \(\mathbf{E}\) 是 Hausdorff 局部凸空间，\(\mathbf{A}\) 是 \(\mathbf{C}\) （\(\mathbf{E}\) 中一个闭凸集）的一个帽。证明：如果 \(a \in \mathbf{A}\) 是 \(\mathbf{A}\) 的一个极点，那么刻面 \(\mathbf{F}\) （\(a\) 在 \(\mathbf{C}\) 中的刻面）（II, 第 87 页, 习题 3）的维数 \(\leqslant 1\) （利用 II, 第 91 页, 习题 26）。由此推出 \(\mathbf{F} \cap \mathbf{A}\) 是 \(\mathbf{C}\) 的一个帽。
\end{enumerate}

¶ * 31) 设 X 是 R 的紧区间 \([0, 1]\) ，\(\Phi\) 是由 X 上定义的连续实值函数与函数 \(t \mapsto |t - a|^{-\alpha}\) （其中 \(a \in X\) 且 \(0 < \alpha < 1\) ；我们置 \(0^{-\alpha} = +\infty\) （对 \(\alpha > 0\) ））所构成的集合。在 X 上测度所成的空间 \(\mathcal{M}(X)\) 中，设 \(M_{\Phi}^{+}\) 是满足如下条件的测度 \(\mu \geqslant 0\) 所成之集：\(\Phi\) 的所有函数都是 \(\mu\) -可积的。

\begin{enumerate}
\def\labelenumi{\alph{enumi})}
\item
  给 \(\mathcal{M}_{\Phi}^{+}\) 装备由 \(\mathbf{R}^{\Phi}\) 的乘积结构诱导的一致结构。证明：对这个一致结构，\(\mathcal{M}_{\Phi}^{+}\) 是一个真的、凸的、完备的锥。（注意：对每个函数 \(f \in \Phi\) ，存在 \(g \in \Phi\) ，使得对所有 \(\varepsilon > 0\) ，存在 \(u \in \mathcal{C}(\mathbf{X}; \mathbf{R})\) 满足 \(0 \leqslant f - u \leqslant \varepsilon g\) 。b) 证明：锥 \(\mathcal{M}_{\Phi}^{+}\) 没有极生成元。（注意：如果 \(\mu \in \mathcal{M}_{\Phi}^{+}\) ，那么所有测度 \(\lambda\) （满足 \(0 \leqslant \lambda \leqslant \mu\) ）都属于 \(\mathcal{M}_{\Phi}^{+}\) 。）
\item
  证明：由 \(\mu \in \mathcal{M}_{\Phi}^{+}\) （满足 \(\mu(1) = 1\) ）所成之集 S 是锥 \(\mathcal{M}_{\Phi}^{+}\) 的一个底，并且是 \(\mathbf{R}^{\Phi}\) 中的一个单纯形（II, 第 71 页, 习题 41）。 *
\end{enumerate}

\begin{enumerate}
\def\labelenumi{\arabic{enumi})}
\setcounter{enumi}{31}
\tightlist
\item
  设 \(\mathbf{E}\) 与 \(\mathbf{F}\) 是两个 Hausdorff 局部凸空间，\(\mathbf{A}\) 是 \(\mathbf{E}\) 的一个凸子集，\(u\) 是 \(\mathbf{E}\) 到 \(\mathbf{F}\) 的一个线性映射。
\end{enumerate}

\begin{enumerate}
\def\labelenumi{\alph{enumi})}
\item
  在 \(u\) 下，\(u(\mathbf{A})\) 的一个支撑流形（II, 第 87 页, 习题 3, c)）的逆像是 \(\mathbf{A}\) 的一个支撑流形。
\item
  如果 \(\mathbf{A}\) 是紧的且 \(u\) 连续，那么 \(u(\mathbf{A})\) 的每个极点都是在 \(u\) 之下 \(\mathbf{A}\) 的一个极点的像。
\item
  如果 \(\mathbf{A}\) 是一个以 0 为顶点的局部紧锥，且 \(u\) 连续，那么 \(u(\mathbf{A})\) 的每个极生成元都是在 \(u\) 之下 \(\mathbf{A}\) 的一个极生成元的像。
\end{enumerate}

¶ 33) 设 E 是一个 Hausdorff 局部凸空间，A 是 E 的一个子集。

\begin{enumerate}
\def\labelenumi{\alph{enumi})}
\item
  用 \(\Gamma_0(\mathbf{A})\) 表示这样的点 \(x \in \mathbf{E}\) 所成之集：对每个连续线性映射 \(u\) （从 \(\mathbf{E}\) 到一个有限维向量空间的映射），像 \(u(x)\) 属于 \(u(\mathbf{A})\) 的凸包。这等价于说：对每个闭线性流形 \(\mathbf{V}\) （\(\mathbf{E}\) 中包含 \(x\) 且余维数为有限数 \(n > 0\) 的闭线性流形），存在 \(\mathbf{A}\) 的一个至多含 \(n + 1\) 个元素的子集，其凸包与 \(\mathbf{V}\) 相交。证明：\(\Gamma_0(\mathbf{A})\) 是一个包含 \(\mathbf{A}\) 的凸集，\(\Gamma_0(\Gamma_0(\mathbf{A})) = \Gamma_0(\mathbf{A})\) ，且 \(\Gamma_0(\mathbf{A})\) 含于 \(\mathbf{A}\) 的闭凸包（利用 II, 第 38 页, 命题 4）。
\item
  设 \((x_{\alpha})_{\alpha \in \mathbb{I}}\) 是 A 的元素的一个族，\((\lambda_{\alpha})_{\alpha \in \mathbb{I}}\) 是正数的一个族，使得 \(\sum_{\alpha \in \mathbb{I}}\lambda_{\alpha} = 1\) 且族 \((\lambda_{\alpha}x_{\alpha})\) 在 E 中可和。证明：和 \(s = \sum_{\alpha \in \mathbb{I}}\lambda_{\alpha}x_{\alpha}\) 属于
\end{enumerate}

\(\Gamma_0(\mathbf{A})\) 。（借助 \(a\) ）化归为 \(\mathbf{E}\) 有限维且等同于由诸 \(\lambda_{\alpha}x_{\alpha}\) 生成的线性流形的情形；然后用反证法论证：对每个有限子集 \(\mathbf{J}\) （\(\mathbf{I}\) 的有限子集），考虑一个闭超平面 \(\mathbf{H}_{\mathbf{J}}\) （它过 \(s\) ，且不与诸 \(x_{\alpha}\) （满足 \(\alpha \in \mathbf{J}\) ）所成之集的凸包相交），再利用有限维空间中单位球的紧性。）

\begin{enumerate}
\def\labelenumi{\alph{enumi})}
\setcounter{enumi}{2}
\item
  证明：如果 A 是紧的，那么 \(\Gamma_0(A)\) 等同于 A 的闭凸包。d) 如果 K 是 E 中一个紧凸集，而 A 是其极点所成之集，证明 \(K = \Gamma_0(A)\) （利用 II, 第 90 页, 习题 22, b) 与习题 32）。
\item
  沿用 II, 第 74 页, 习题 3 的记号，设 A 是由诸 \(\varepsilon_{x}\) 构成的集合，其中 \(x\) 取遍满足 \(0 \leqslant x \leqslant 1\) 的有理数。证明：\(\Gamma_0(A)\) 异于 A 的凸包，也异于 A 的闭凸包。
\end{enumerate}

¶ 34) 设 S 是 E（一个 Hausdorff 局部凸空间）中一个闭凸集，A 是 S 的一个子集，使得 \(\mathbf{S} = \Gamma_0(\mathbf{A})\) （习题 33），并设 \(\mathbf{S}_0\) 是 A 的凸包（于是 \(\mathbf{S} = \bar{\mathbf{S}}_0\) ）。设 N 是 E 的一个闭凸子集，它包含一个余维数为有限的闭线性流形；令 \(\mathbf{M} = \mathbf{S} \cap \mathbf{N}\) ，\(\mathbf{M}_0 = \mathbf{S}_0 \cap \mathbf{N}\) 。

\begin{enumerate}
\def\labelenumi{\alph{enumi})}
\item
  证明：\(\mathbf{M} = \overline{\mathbf{M}}_0\) 。（利用习题 33, a) 注意，E 的每个包含一点 \(x \in \mathbf{M}\) 且余维数为有限的闭线性流形都与 \(\mathbf{M}_0\) 相交，并利用 II, 第 38 页, 命题 4。）
\item
  假设对 A 的每个有限子集 F，N 与 F 的凸包的每个点在 S 中的刻面之交是紧的或有限维的且不含直线。证明：M 是其极点所成之集的闭凸包。（借助 a)，这化归为证明 \(M_{0}\) 的每一点都含于 M 的极点所成之集的闭凸包。利用 II, 第 87 页, 习题 3, e)、Krein--Milman 定理（II, 第 55 页）与 II, 第 90 页, 习题 22, b)。）由此推出：M 的每个闭支撑超平面都含 M 的一个极点。
\end{enumerate}

\begin{enumerate}
\def\labelenumi{\arabic{enumi})}
\setcounter{enumi}{34}
\tightlist
\item
  设 I 是一个不可数集，记 \(E = R^{(I)}\) ，\(E' = E \times R\) 。用 \((e_{\alpha})_{\alpha \in I}\) 表示 E 的典范基，并设 s 是 \((0, 1)\) （\(E'\) 的元素）。在 E 与 \(E'\) 之间定义一个分离对偶，即令 \(\langle e_{\alpha}, e_{\beta} \rangle = \delta_{\alpha\beta}, \langle e_{\alpha}, s \rangle = 1\) （对所有 \(\alpha \in I\) ）。设 C 是 E 中的尖锥 \(R_{+}^{I}\) 。a) 证明：由 \(\sigma(E, E')\) 与 E 上的范数 \(p(x) = \sum_{\alpha \in I} |x_{\alpha}|\) 在 C 上诱导的拓扑重合。
\end{enumerate}

\begin{enumerate}
\def\labelenumi{\alph{enumi})}
\setcounter{enumi}{1}
\tightlist
\item
  证明：由 \(\sigma(\mathrm{E},\mathrm{E}^{\prime})\) 在 C 上诱导的一致结构不是可度量化的。
\end{enumerate}

\begin{enumerate}
\def\labelenumi{\arabic{enumi})}
\setcounter{enumi}{35}
\tightlist
\item
  考虑空间 \(\mathrm{E} = \mathbf{R}^{(\mathbf{N})}\) ，带弱拓扑 \(\sigma (\mathbf{R}^{(\mathbf{N})},\mathbf{R}^{\mathbf{N}})\) ；设 C 是 E 中由点 \(x = (x_{n})\) （满足 \(x_{n}\geqslant 0\) ，对所有 \(n\) ）所构成的闭凸锥。\\
\end{enumerate}

\begin{enumerate}
\def\labelenumi{\alph{enumi})}
\item
  设 \(x = (x_{n})\) 是 C 的一点，J 是整数 \(n\) （满足 \(x_{n} > 0\) ）所成的有限集；如果 \(m\) 是 J 的元素个数，那么设 A 是 C 中这样的点 \(y = (y_{n})\) 所成之集：\(y_{n} = 0\) （对 \(n\in \mathrm{J}\) ）且 \(\sum_{k\in J}y_kx_k^{-1}\leqslant m\) 。证明：A 是 C 的一个包含 \(x\) 的帽。
\item
  证明：不存在 C 中的帽 B，使得 C 是诸集合 nB（n \textgreater{} 0）之并。（设 p 是 B 的度规在 C 上的限制；p 在 C 上有限，且如果 \((e_{n})\) 是 E 的典范基，那么 \(p(e_{n}) > 0\) 对所有 n 成立（II, 第 58 页, 命题 4），而诸点 \(z^{(n)} = e_{n}/p(e_{n})\) 属于 B；但证明：存在 \(z' \in R^{N}\) ，使序列 \((\langle z^{(n)}, z' \rangle)\) 不是有界的。）
\end{enumerate}

\begin{enumerate}
\def\labelenumi{\arabic{enumi})}
\setcounter{enumi}{36}
\tightlist
\item
  设 F 是 Banach 空间 \(l^{1}(\mathbf{N})\) ，它由实数的可和序列 \(x = (x_{n})\) 组成；E 是趋于 0 的序列 \(y = (y_{n})\) 所成的空间；我们给 F 装备弱拓扑 \(\sigma(\mathrm{F}, \mathrm{E})\) ，这里 E 与 F 借助型 \(\mathbf{B}(x, y) = \sum x_{n} y_{n}\) 处于分离对偶。
\end{enumerate}

\begin{enumerate}
\def\labelenumi{\alph{enumi})}
\item
  设 \(\mathbf{C}\) 是 \(\mathbf{F}\) 中由点 \(x = (x_{n})\) （满足 \(x_{n} \geqslant 0\) ，对所有 \(n\) ）所构成的凸锥。证明：\(\mathbf{C}\) 在 \(\mathbf{F}\) 中是闭的。
\item
  设 A 是由 \(x = (x_{n}) \in \mathbf{C}\) （满足 \(\sum_{n} x_{n} \leqslant 1\) ）所成之集。证明：A 是 C 的一个帽，它对由 F 的拓扑诱导的拓扑是可度量化的，且 C 是诸集合 nA（\(n > 0\) ）之并。
\item
  证明：C 的底不是紧的（这样一个集合 S 将是由 \(x = (x_{n}) \in \mathbf{C}\) （满足 \(\sum_{n} z_{n} x_{n} = 1\) ，其中 \((z_{n}) \in \mathbf{E}\) 且 \(z_{n} > 0\) 对所有 n）所成之集。如果 \(e_{m} = (\delta_{mn})_{n \geqslant 0}\) ，诸点 \(z_{n}^{-1} e_{n}\) 属于 S，但不构成 F 中的相对紧集）。
\end{enumerate}

\(d\) ) 证明：\(\mathbf{C}\) 对由 \(\mathbf{F}\) 的拓扑诱导的拓扑不是可度量化的（利用 Baire 定理，并注意 A 中没有一点是相对于子空间 C 的代数内点）。

\begin{enumerate}
\def\labelenumi{\arabic{enumi})}
\setcounter{enumi}{37}
\tightlist
\item
  设 E 是一个 Hausdorff 局部凸空间，X 是 E 中一个紧凸集。用 \(\mathcal{A}(X)\) 表示 X 上连续仿射函数所成之集（不必是 E 上连续仿射函数在 X 上的限制，参见 II, 第 78 页, 习题 11）。对 X 上有上界的每个实值函数 f，令 \(\tilde{f}(x) = \inf_{h}(h(x))\) （对所有 \(x \in X\) ），其中 h 取遍 \(\mathcal{A}(X)\) 中满足 \(h \geqslant f\) 的函数所成之集。
\end{enumerate}

\begin{enumerate}
\def\labelenumi{\alph{enumi})}
\item
  证明：\(\tilde{f}\) 是上半连续的凹函数。如果 \(f\) 本身上半连续且凹，那么 \(\tilde{f} = f\) （参见 II, 第 39 页, 命题 5）。
\item
  假设 \(f\) 上半连续。证明：\(\tilde{f}\) 是诸函数 \(\tilde{g}\) （\(g\) 取遍一个函数集合，其中的函数在 \(X\) 上连续且以 \(f\) 为下包络）的下包络。
\item
  对所有 \(x \in \mathbf{X}\) ，用 \(\mathcal{M}_x\) 表示有限族 \(\mu = ((\mu_j, x_j))\) 所成之集，其中 \(x_j\) 是 \(\mathbf{X}\) 的点，诸 \(\mu_j\) 是数 \(\geqslant 0\) ，满足 \(\sum_j \mu_j = 1\) ，使得 \(x = \sum_j \mu_j x_j\) 。对每个实值函数 \(f\) （在 \(\mathbf{X}\) 上有上界），令 \(f'(x) = \sup_{\mu \in \mathcal{M}_x} \sum \mu_j f(x_j)\) （对所有 \(x \in \mathbf{X}\) ）。证明：\(f'\) 是 \(\mathbf{X}\) 上的一个凹函数，且 \(f' \leqslant \tilde{f}\) 。
\item
  假设 \(f\) 在 \(\mathbf{X}\) 上连续。给定 \(\varepsilon > 0\) ，设 \((\mathrm{U}_k)_{1 \leqslant k \leqslant N}\) 是 \(\mathbf{X}\) 的一个由凸开集组成的覆盖，使得关系 \(x \in \mathrm{U}_k\) 与 \(y \in \mathrm{U}_k\) 蕴含不等式 \(|f(x) - f(y)| \leqslant \varepsilon\) 。令 \(\mathrm{A}_1 = \mathrm{U}_1\) ，且对 \(k > 1\) ，令 \(\mathrm{A}_k = \mathrm{U}_k \cap \complement(\mathrm{U}_1 \cup \mathrm{U}_2 \cup \mathrm{U}_3 \cup \ldots \cup \mathrm{U}_{k-1})\) 。证明：对所有 \(x \in \mathbf{X}\) ，存在一个族 \(\mu = ((\mu_k, x_k))\) （由 \(N\) 项组成，属于 \(\mathcal{M}_x\) ），其中 \(x_k \in \mathrm{U}_k\) （对 \(1 \leqslant k \leqslant N\) ），使得 \(\sum_k \mu_k f(x_k) \geqslant f'(x) - 2\varepsilon\) （如果我们有不等式 \(\sum_j \lambda_j f(y_j) \geqslant f'(x) - \varepsilon\) ，就把诸 \(y_j\) （属于同一个 \(A_k\) 的那些）归并在一起）。
\item
  由 \(d\) 推出：当 \(f\) 连续时，\(f'\) 上半连续且 \(f' = \tilde{f}\) 。（如果 \(\mathfrak{U}\) 是 \(\mathbf{X}\) 上一个比某点 \(x \in \mathbf{X}\) 的邻域滤子更细的超滤子，且若 \(f'(y) \geqslant r\) 对所有这样的点 \(y\) 成立：这些点属于 \(\mathfrak{U}\) 中的某个集合，那么证明 \(f'(x) \geqslant r - 2\varepsilon\) ，办法是让一个族 \(\mu_y \in \mathcal{M}_y\) 对应于每个 \(y\) （该族满足条件 \(d\) ) ），然后沿 \(\mathfrak{U}\) 取极限。）
\end{enumerate}

\begin{enumerate}
\def\labelenumi{\arabic{enumi})}
\setcounter{enumi}{38}
\tightlist
\item
  设 \(\mathbf{H}\) 是 Hausdorff 局部凸空间 \(\mathbf{E}\) 中一个不包含 0 的闭超平面，\(S\) 是含于 \(\mathbf{H}\) 中的一个紧单纯形（II, 第 71 页, 习题 41）。
\end{enumerate}

\begin{enumerate}
\def\labelenumi{\alph{enumi})}
\item
  设 C 是由 S 生成的以 0 为顶点的锥。证明：如果 \((x_{i})_{i\in \mathbf{I}}\) 与 \((y_{j})_{j\in \mathbf{J}}\) 是 C 的点的两个有限族，满足 \(\sum_{i\in \mathbf{I}}x_i = \sum_{j\in \mathbf{J}}y_j\) ，那么存在 C 的点的一个有限族 \((z_{ij})_{(i,j)\in \mathbf{I}\times \mathbf{J}}\) ，使得 \(x_{i} = \sum_{j\in \mathbf{J}}z_{ij}\) 对所有 \(i\in \mathbf{I}\) 成立，且 \(y_{j} = \sum_{i\in \mathbf{I}}z_{ij}\) 对所有 \(j\in \mathbf{J}\) 成立（用归纳法论证，化归到 \(\mathrm{I} = \mathrm{J} = \{1,2\}\) 的情形）。
\item
  设 \(f\) 是 S 上一个上半连续的凸函数。证明：函数 \(\tilde{f}\) （在习题 38 中定义）是一个仿射函数。（首先利用习题 38, b) 与 II, 第 81 页, 习题 28 化归到 \(f\) 连续的情形。然后利用这样的事实：如果 \(f\) 连续，那么 \(\tilde{f} = f'\) （习题 38, e)），并证明 \(f'\) 是凸的，办法是利用 a) 对 \(f'(\alpha_1 x_1 + \alpha_2 x_2)\) 加以上界估计，其中 \(\alpha_1 \geqslant 0\) ，\(\alpha_2 \geqslant 0\) ，\(\alpha_1 + \alpha_2 = 1\) 。）
\end{enumerate}

\begin{enumerate}
\def\labelenumi{\arabic{enumi})}
\setcounter{enumi}{39}
\tightlist
\item
  设 \(\mathbf{X}\) 是 \(\mathbf{E}\) （一个 Hausdorff 局部凸空间）中一个紧凸集，\(f\) 是一个有下界的上半连续函数。设 \(g\) 是一个下半连续的凹函数，满足 \(g \geqslant f\) 。证明（用习题 38 的记号）：\(g \geqslant \tilde{f}\) 。（化归到 \(\inf_{x \in \mathbf{X}} (g(x) - f(x)) > 0\) 的情形。如果 \((f_{\alpha})\) 是由连续
\end{enumerate}

函数组成的一个递减有向族，使得 \(f = \inf(f_{\alpha})\) ，证明此时也有 \(\inf_{x \in X} (g(x) - f_{\alpha}(x)) \geqslant 0\) 对 \(\alpha \geqslant \alpha_0\) 成立，从而把问题化归到 \(f\) 连续的情形。然后利用习题 38, e)。

\begin{enumerate}
\def\labelenumi{\arabic{enumi})}
\setcounter{enumi}{40}
\tightlist
\item
  设 S 是含于一个 Hausdorff 局部凸空间 E 中的一个紧单纯形（II, 第 71 页, 习题 41），\(f\) 是一个有下界的上半连续凸函数。设 \(g\) 是一个下半连续的凹函数，满足 \(g \geqslant f\) 。
\end{enumerate}

\begin{enumerate}
\def\labelenumi{\alph{enumi})}
\item
  如果记 \(u = \tilde{f}\) ，\(v = -(-g)\tilde{f}\) ，那么 \(u\) 与 \(v\) 是满足 \(u \leqslant v\) 的仿射函数（利用习题 39, b) 与习题 40）。
\item
  证明：存在一个仿射函数 \(h\) ，它在 \(\mathbf{X}\) 上连续且满足 \(f \leqslant h \leqslant g\) （D. Edwards 定理）。（我们可以把 \(f\) 换成 \(u\) ，把 \(g\) 换成 \(v\) 。构造仿射函数的三个序列 \((u_m), (v_m), (h_m)\) ，使得在 \(\mathbf{X}\) 上，函数 \(u_m\) 上半连续，函数 \(v_m\) 下半连续，函数 \(h_m\) 连续，且
\end{enumerate}

\[
u - \frac {1}{2 ^ {m}} \leqslant u _ {m} <   h _ {m} <   v _ {m} \leqslant v + \frac {1}{2 ^ {m}} \text { and } \| h _ {n + 1} - h _ {n} \| \leqslant \frac {1}{2 ^ {n + 1}}.
\]

为此利用 II, 第 81 页, 习题 29。）

